\chapter{KẾT LUẬN VÀ HƯỚNG PHÁT TRIỂN}

\section{Kết luận}

Đồ án đã xây dựng thành công website đặt lịch cắt tóc trực tuyến Barber
Shop với đầy đủ các chức năng cốt lõi của một hệ thống quản lý dịch vụ.
Hệ thống đáp ứng tương đối đầy đủ ba nhóm người dùng chính gồm khách
hàng, thợ cắt tóc và quản trị viên, từ đó mô phỏng được luồng nghiệp vụ
thực tế của một tiệm cắt tóc hiện đại.

Đối với khách hàng, hệ thống hỗ trợ đăng ký, đăng nhập, đăng nhập bằng
Google, xem dịch vụ, chọn chi nhánh, chọn thợ, đặt lịch, kiểm tra lịch
sử dụng, lịch sử đặt hẹn và gửi đánh giá sau khi hoàn thành dịch vụ.
Đối với thợ cắt tóc, hệ thống cho phép theo dõi lịch làm việc và gửi đơn
xin nghỉ phép. Đối với quản trị viên, hệ thống cung cấp các chức năng
quản lý khách hàng, thợ, dịch vụ, chi nhánh, lịch hẹn, đơn nghỉ phép,
đánh giá và thống kê tổng quan.

Trong quá trình triển khai, đồ án đã chú trọng đến tính rõ ràng của giao
diện, tính nhất quán trong xử lý dữ liệu và khả năng mở rộng của hệ
thống. Việc tổ chức theo mô hình Client -- Server với PHP, MySQL, HTML,
CSS và JavaScript giúp tách biệt tương đối rõ giữa giao diện và xử lý
nghiệp vụ, đồng thời tạo thuận lợi cho việc bảo trì và phát triển về sau.

\section{Các hạn chế cần khắc phục}

Mặc dù hệ thống đã hoàn thiện các chức năng nghiệp vụ chính, vẫn còn một
số hạn chế có thể tiếp tục cải thiện. Trước hết, phần thống kê mới dừng
ở mức cung cấp số liệu tổng quan, chưa triển khai đầy đủ các biểu đồ
phân tích chuyên sâu theo nhiều chiều như hiệu suất từng thợ theo tuần,
so sánh doanh thu theo chi nhánh hoặc dự báo xu hướng đặt lịch.

Bên cạnh đó, hệ thống chưa tích hợp cơ chế thông báo tự động qua email
hoặc tin nhắn để nhắc lịch trước thời gian hẹn. Người dùng vẫn cần chủ
động truy cập website để kiểm tra trạng thái lịch hẹn, vì vậy trải
nghiệm nhắc lịch chưa thật sự tối ưu.

Ngoài ra, việc quản trị hệ thống vẫn có thể được nâng cấp thêm ở các
khía cạnh như phân quyền chi tiết hơn cho từng nhóm người dùng, ghi log
thao tác quản trị, và chuẩn hóa quy trình kiểm duyệt dữ liệu khi hệ
thống phát sinh nhiều giao dịch đồng thời. Đây là những hướng cần thiết
để hệ thống vận hành ổn định hơn trong môi trường thực tế.

\section{Hướng phát triển trong tương lai}

Dựa trên những hạn chế đã được nhìn nhận, các hướng phát triển trong
tương lai sẽ tập trung vào việc nâng cao mức độ tự động hóa và khả năng
hỗ trợ quản trị. Trước hết, hệ thống có thể được mở rộng thêm các biểu
đồ thống kê nâng cao, dashboard tùy biến và báo cáo định kỳ để giúp
quản trị viên theo dõi hiệu quả kinh doanh theo từng giai đoạn.

Tiếp theo, việc tích hợp thông báo tự động qua email hoặc Zalo/SMS sẽ
giúp khách hàng nhận xác nhận lịch hẹn, lời nhắc trước giờ phục vụ và
thông tin thay đổi trạng thái kịp thời hơn. Điều này không chỉ cải thiện
trải nghiệm người dùng mà còn giảm thiểu trường hợp quên lịch hoặc đến
sai giờ.

Ngoài ra, hệ thống có thể tiếp tục phát triển các chức năng như thanh
toán trực tuyến, mã khuyến mãi, quản lý vật tư, lịch làm việc nâng cao
cho thợ và cơ chế phân tích đánh giá khách hàng. Nếu được đầu tư thêm
về dữ liệu và giao diện, hệ thống sẽ ngày càng hoàn thiện hơn, phù hợp
hơn với mô hình vận hành thực tế của một tiệm cắt tóc chuyên nghiệp.

